\chapter{Engineering as a way of reasoning}

\section{The object of the discipline}
Programming is the act of expressing a computation. \term{Software engineering}
adds a harder question: how can that computation remain correct, understandable,
secure, and economically maintainable while its environment changes? The
environment includes users, regulations, dependencies, hardware, networks,
other teams, and the people who will inherit the code without the original
author's memory.

The word \emph{engineering} is useful because it points to constrained design.
An engineer makes assumptions explicit, models important failure modes, chooses
materials and boundaries, measures the result, and manages uncertainty. Software
has unusual properties: it is easy to copy, cheap to modify in one place, and
often difficult to observe as a whole. A change that looks local can alter
security, latency, data interpretation, or the behaviour of a distant client.

\begin{keyidea}
Quality is not a single attribute. It is the degree to which a system satisfies
the requirements that matter in its context, including requirements about
change, operation, safety, privacy, and future understanding.
\end{keyidea}

\section{Claims, evidence, and recommendations}
Technical writing often mixes three kinds of statement. A specification defines
what an implementation must do. Empirical research reports observations under
specified conditions. Professional practice is accumulated experience that may
be useful without being a universal law. A recommendation is a judgment about a
particular context. Keeping these categories separate prevents cargo cults.

For example, HTTP semantics are defined by an IETF standard; the assertion that
an HTTP method is safe or idempotent is a protocol claim that can be checked
against that standard \cite{ietf2022http}. The belief that small pull requests
are easier to review is a practice hypothesis: often plausible, but dependent
on change shape, reviewer skill, and system complexity. The decision to use a
modular monolith for a small product is an engineering recommendation whose
strength comes from the constraints, not from the label.

\section{A system boundary is a decision}
Every system description has a boundary. Inside it are behaviours the team can
change directly; outside it are actors, services, devices, laws, and physical
conditions. A boundary is never purely technical. It encodes ownership and
responsibility.

The Ledgerline boundary is deliberately modest:

\begin{itemize}
  \item a caller submits an order containing a customer identifier and line items;
  \item the service validates domain invariants and returns a stable order identifier;
  \item a relational store records the accepted state;
  \item operators can tell whether the service is alive, useful, and overloaded;
  \item a failure does not silently create a second order when the caller retries.
\end{itemize}

The last point is not an implementation detail. It is a business property that
leads to an idempotency key, a database constraint, a response contract, and a
recovery test. Engineering begins when a vague concern becomes a checkable
property.

\section{The feedback loop}
Good development is a loop rather than a waterfall of documents:

\begin{center}
\begin{tikzpicture}[node distance=6mm,>=Latex,font=\small]
  \node[draw,rounded corners,fill=codebg,minimum width=22mm,minimum height=8mm] (need) {need};
  \node[draw,rounded corners,fill=codebg,minimum width=22mm,minimum height=8mm,right=of need] (model) {model};
  \node[draw,rounded corners,fill=codebg,minimum width=22mm,minimum height=8mm,right=of model] (build) {build};
  \node[draw,rounded corners,fill=codebg,minimum width=22mm,minimum height=8mm,right=of build] (observe) {observe};
  \draw[->,thick,accent] (need) -- (model);
  \draw[->,thick,accent] (model) -- (build);
  \draw[->,thick,accent] (build) -- (observe);
  \draw[->,thick,accent] (observe) to[bend left=35] (need);
\end{tikzpicture}
\end{center}

The model may be a test, a state machine, a schema, a sequence diagram, or a
small written invariant. The point is not to predict every detail. It is to
make the important assumptions cheap to challenge before they become expensive
code or production incidents.

\section{A practical quality model}
Before choosing tools, ask which of these dimensions are material:

\noindent\begin{tabularx}{\textwidth}{@{}p{0.22\textwidth}X@{}}
\toprule
Dimension & Question to answer \\
\midrule
Correctness & Does the system produce the required result for valid and invalid inputs? \\
Safety and security & What must never happen, and who might deliberately cause it? \\
Reliability & What happens when dependencies, data, or operators behave unexpectedly? \\
Operability & Can a responsible person detect, diagnose, and recover from failure? \\
Changeability & Can a team modify one concern without guessing about unrelated consequences? \\
Human fit & Can people with different abilities, languages, roles, and incentives use and govern it? \\
\bottomrule
\end{tabularx}

These dimensions compete. Stronger durability may cost latency; a highly generic
API may be harder to understand; a fast migration may increase rollback risk.
The engineering response is not to maximize every quality. It is to state the
trade-off and make the consequence visible.

\begin{exercise}
Write five falsifiable properties for a service that accepts payments. For each,
name the boundary at which it is measured and one test or observation that could
show it is false.
\end{exercise}

\section{Chapter summary}
Software engineering is disciplined reasoning under constraints. Begin with
properties rather than technologies, distinguish standards from evidence and
recommendations, draw a deliberate system boundary, and use short feedback
loops to turn assumptions into tests and observations.
